\section{Related Work}
\label{sec:related}

\textbf{Capability-focused safety benchmarks.} WMDP~\citep{li2024wmdp} measures whether models possess dangerous knowledge across biosecurity, cybersecurity, and chemical weapons domains. RAND's evaluation of AI-assisted biological threat capabilities~\citep{mouton2024operational} assesses whether AI assistance provides meaningful uplift. HarmBench~\citep{mazeika2024harmbench} provides a standardized framework for evaluating attack and defense methods across diverse harm categories. SafetyBench~\citep{zhang2024safetybench} offers a comprehensive multiple-choice benchmark for safety evaluation across 11 categories in both English and Chinese. These benchmarks answer ``what does the model know?'' or ``can it be attacked?'' but not ``do mitigations work under realistic conditions?''

\textbf{Uplift measurement.} METR's AI-biology randomized controlled trial~\citep{metr2026} provides the gold standard for measuring whether AI assistance changes experimental outcomes. However, RCTs are expensive, domain-specific, and infeasible for continuous monitoring of mitigation effectiveness across model versions.

\textbf{Safety frameworks.} DeepMind's Frontier Safety Framework~\citep{deepmind2024fsf} defines Critical Capability Levels (CCLs) for evaluating model risk. Anthropic's Responsible Scaling Policy~\citep{anthropic2023rsp} introduces AI Safety Levels (ASLs) with corresponding evaluation commitments. OpenAI's Preparedness Framework~\citep{openai2023preparedness} provides risk categorization. These frameworks motivate evaluation but do not provide reusable, contamination-aware evaluation toolkits.

\textbf{LLM-as-judge evaluation.} MT-Bench and Chatbot Arena~\citep{zheng2023judging} demonstrate that strong LLMs can serve as reliable evaluators, achieving high agreement with human preferences. G-Eval~\citep{liu2023geval} formalizes LLM-based evaluation with chain-of-thought reasoning. We adopt LLM-as-judge scoring as a validation layer for our pattern-based scorer, enabling dual-scoring inter-rater analysis without requiring human annotator recruitment.

\textbf{Evaluation methodology.} Recent work on benchmark saturation~\citep{kiela2021dynabench}, contamination detection~\citep{oren2023proving}, and holistic evaluation~\citep{liang2023holistic} highlights the fragility of static benchmarks. Goodhart's law in AI evaluation~\citep{manheim2019categorizing} warns that metrics cease to be useful when they become targets. Our work applies these insights specifically to safety-sensitive evaluation, where contaminated or gamed benchmarks can create false assurance about model safety.

\textbf{Position.} We complement capability benchmarks (which measure knowledge) and attack benchmarks (which measure adversarial robustness) with a methodology for testing whether \emph{deployment mitigations} actually hold on sensitive scientific content, using contamination-aware design and governed release infrastructure.
